% GENERATED FILE. Edit canonical lessons, then run npm run generate:books.
% canonical-source-hash: fnv1a64:08f242cf95b15489
% canonical-lessons: AR-C450-baladiyya, AR-C450-qunsuliyya, AR-C450-dawam, AR-C450-tabur, AR-C450-taqaddama

\chapter{Town hall and consulate}
\label{ch:ar-town-hall-and-consulate}
\hlchaptermodality{450}

\begin{chapteropening}
\textbf{By the end of this chapter:} \emph{I can say a town hall, a consulate, working hours, a line of people and to apply.}

Complete the last lesson of chapter 450: I can say a town hall, a consulate, working hours, a line of people and to apply.
\end{chapteropening}

\section[\texorpdfstring{\ar{بلدية} (baladiyya)}{بلدية (baladiyya)}]{\ar{بلدية} (baladiyya) — a town hall; a municipality}

\label{lesson:AR-C450-baladiyya}



\begin{quote}
Before the new one: say the Arabic for to grumble, then the Arabic for to swap.
\end{quote}

\subsection*{You'll want to know: \ar{بلدية}}
\textbf{\ar{بلدية}} — \emph{baladiyya} — \textquotedblleft{}a town hall; a municipality\textquotedblright{}.

\begin{quote}
\ar{ذهبت} \ar{إلى} \ar{البلدية} \ar{لدفع} \ar{الرسوم}.
\end{quote}

\begin{quote}
\emph{dhahabtu ilā al-baladiyya li-dafʿ ar-rusūm.}
\end{quote}

\begin{quote}
I went to the town hall to pay the fees.
\end{quote}

\begin{cousinweb}
\textbf{\ar{بلدية}} (\emph{baladiyya}) is built on the native Arabic root \emph{b-l-d}, \textquotedblleft{}a town, a country\textquotedblright{} (Lane). It is a noun in \emph{-iyya} made from \emph{balad}: the body that runs a town.
\end{cousinweb}

\begin{grammarlens}[title={What you've built}]
One more word for getting around.
\end{grammarlens}

\subsection*{Your turn}
\begin{itemize}
\raggedright
  \item \emph{Say it:} \emph{baladiyya}
  \item \emph{Say it:} \emph{baladiyya}, once more
  \item \emph{Say it:} \emph{dhahabtu ilā al-baladiyya li-dafʿ ar-rusūm.} — I went to the town hall to pay the fees.
  \item \emph{Recall:} say the Arabic for to grumble, then the Arabic for to swap, then say \emph{baladiyya} again
\end{itemize}

\begin{tcolorbox}[breakable,colback=teal!4,colframe=teal!35!black,title={Before you move on}]
What does \ar{بلدية} mean? (\textquotedblleft{}A town hall; a municipality\textquotedblright{}.) Say it once more.
\end{tcolorbox}
\section[\texorpdfstring{\ar{قنصلية} (qunṣuliyya)}{قنصلية (qunṣuliyya)}]{\ar{قنصلية} (qunṣuliyya) — a consulate}

\label{lesson:AR-C450-qunsuliyya}



\begin{quote}
Before the new one: say the Arabic for to swap, then the Arabic for a town hall.
\end{quote}

\subsection*{You'll want to know: \ar{قنصلية}}
\textbf{\ar{قنصلية}} — \emph{qunṣuliyya} — \textquotedblleft{}a consulate\textquotedblright{}.

\begin{quote}
\ar{القنصلية} \ar{مغلقة} \ar{يوم} \ar{الجمعة}.
\end{quote}

\begin{quote}
\emph{al-qunṣuliyya mughlaqa yawm al-jumʿa.}
\end{quote}

\begin{quote}
The consulate is closed on Friday.
\end{quote}

\begin{cousinweb}
\textbf{\ar{قنصلية}} (\emph{qunṣuliyya}) is usually traced to the European word \emph{consul}, from Latin, with the Arabic ending \emph{-iyya}.
\end{cousinweb}

\begin{grammarlens}[title={What you've built}]
One more word for getting around.
\end{grammarlens}

\subsection*{Your turn}
\begin{itemize}
\raggedright
  \item \emph{Say it:} \emph{qunṣuliyya}
  \item \emph{Say it:} \emph{qunṣuliyya}, once more
  \item \emph{Say it:} \emph{al-qunṣuliyya mughlaqa yawm al-jumʿa.} — The consulate is closed on Friday.
  \item \emph{Recall:} say the Arabic for to swap, then the Arabic for a town hall, then say \emph{qunṣuliyya} again
\end{itemize}

\begin{tcolorbox}[breakable,colback=teal!4,colframe=teal!35!black,title={Before you move on}]
What does \ar{قنصلية} mean? (\textquotedblleft{}A consulate\textquotedblright{}.) Say it once more.
\end{tcolorbox}
\section[\texorpdfstring{\ar{دوام} (dawām)}{دوام (dawām)}]{\ar{دوام} (dawām) — working hours; office hours}

\label{lesson:AR-C450-dawam}



\begin{quote}
Before the new one: say the Arabic for a town hall, then the Arabic for a consulate.
\end{quote}

\subsection*{You'll want to know: \ar{دوام}}
\textbf{\ar{دوام}} — \emph{dawām} — \textquotedblleft{}working hours; office hours\textquotedblright{}.

\begin{quote}
\ar{ينتهي} \ar{الدوام} \ar{بعد} \ar{الظهر}.
\end{quote}

\begin{quote}
\emph{yantahī ad-dawām baʿda aẓ-ẓuhr.}
\end{quote}

\begin{quote}
Office hours end in the afternoon.
\end{quote}

\begin{cousinweb}
\textbf{\ar{دوام}} (\emph{dawām}) is built on the native Arabic root \emph{d-w-m}, \textquotedblleft{}to last, to continue\textquotedblright{} (Lane). Working hours are the time one keeps at a job.
\end{cousinweb}

\begin{grammarlens}[title={What you've built}]
One more word for getting around.
\end{grammarlens}

\subsection*{Your turn}
\begin{itemize}
\raggedright
  \item \emph{Say it:} \emph{dawām}
  \item \emph{Say it:} \emph{dawām}, once more
  \item \emph{Say it:} \emph{yantahī ad-dawām baʿda aẓ-ẓuhr.} — Office hours end in the afternoon.
  \item \emph{Recall:} say the Arabic for a town hall, then the Arabic for a consulate, then say \emph{dawām} again
\end{itemize}

\begin{tcolorbox}[breakable,colback=teal!4,colframe=teal!35!black,title={Before you move on}]
What does \ar{دوام} mean? (\textquotedblleft{}Working hours; office hours\textquotedblright{}.) Say it once more.
\end{tcolorbox}
\section[\texorpdfstring{\ar{طابور} (ṭābūr)}{طابور (ṭābūr)}]{\ar{طابور} (ṭābūr) — a line of people; a queue}

\label{lesson:AR-C450-tabur}



\begin{quote}
Before the new one: say the Arabic for a consulate, then the Arabic for working hours.
\end{quote}

\subsection*{You'll want to know: \ar{طابور}}
\textbf{\ar{طابور}} — \emph{ṭābūr} — \textquotedblleft{}a line of people; a queue\textquotedblright{}.

\begin{quote}
\ar{وقفنا} \ar{في} \ar{طابور} \ar{طويل} \ar{أمام} \ar{القنصلية}.
\end{quote}

\begin{quote}
\emph{waqafnā fī ṭābūr ṭawīl amāma al-qunṣuliyya.}
\end{quote}

\begin{quote}
We stood in a long queue in front of the consulate.
\end{quote}

\begin{cousinweb}
\textbf{\ar{طابور}} (\emph{ṭābūr}) was borrowed from Ottoman Turkish (EtymDB), where \emph{tabur} names a battalion, a line of soldiers.
\end{cousinweb}

\begin{grammarlens}[title={What you've built}]
One more word for getting around.
\end{grammarlens}

\subsection*{Your turn}
\begin{itemize}
\raggedright
  \item \emph{Say it:} \emph{ṭābūr}
  \item \emph{Say it:} \emph{ṭābūr}, once more
  \item \emph{Say it:} \emph{waqafnā fī ṭābūr ṭawīl amāma al-qunṣuliyya.} — We stood in a long queue in front of the consulate.
  \item \emph{Recall:} say the Arabic for a consulate, then the Arabic for working hours, then say \emph{ṭābūr} again
\end{itemize}

\begin{tcolorbox}[breakable,colback=teal!4,colframe=teal!35!black,title={Before you move on}]
What does \ar{طابور} mean? (\textquotedblleft{}A line of people; a queue\textquotedblright{}.) Say it once more.
\end{tcolorbox}
\section[\texorpdfstring{\ar{تقدم} (taqaddama)}{تقدم (taqaddama)}]{\ar{تقدم} (taqaddama) — to apply (for); to step forward}

\label{lesson:AR-C450-taqaddama}



\begin{quote}
Before the new one: say the Arabic for working hours, then the Arabic for a line of people.
\end{quote}

\subsection*{You'll want to know: \ar{تقدم}}
\textbf{\ar{تقدم}} — \emph{taqaddama} — \textquotedblleft{}to apply (for); to step forward\textquotedblright{}.

\begin{quote}
\ar{تقدمت} \ar{أختي} \ar{بطلب} \ar{تأشيرة}.
\end{quote}

\begin{quote}
\emph{taqaddamat ukhtī bi-ṭalab taʾshīra.}
\end{quote}

\begin{quote}
My sister applied for a visa.
\end{quote}

Present: \ar{يتقدم}, \emph{yataqaddam}.

\emph{taqaddama bi-ṭalab}: to put in an application.

\begin{cousinweb}
\textbf{\ar{تقدم}} (\emph{taqaddama}) is built on the native Arabic root \emph{q-d-m}, \textquotedblleft{}to come forward, to precede\textquotedblright{} (Lane). It is Form V.
\end{cousinweb}

\begin{grammarlens}[title={What you've built}]
One more word for getting around.
\end{grammarlens}

\subsection*{Your turn}
\begin{itemize}
\raggedright
  \item \emph{Say it:} \emph{taqaddama}
  \item \emph{Say it:} \emph{taqaddama}, once more
  \item \emph{Say it:} \emph{taqaddamat ukhtī bi-ṭalab taʾshīra.} — My sister applied for a visa.
  \item \emph{Recall:} say the Arabic for working hours, then the Arabic for a line of people, then say \emph{taqaddama} again
\end{itemize}

\begin{tcolorbox}[breakable,colback=teal!4,colframe=teal!35!black,title={Before you move on}]
What does \ar{تقدم} mean? (\textquotedblleft{}To apply (for); to step forward\textquotedblright{}.) Say it once more.
\end{tcolorbox}
